\documentclass[11pt]{article}

\usepackage[margin=1in]{geometry}
\usepackage{amsmath,amssymb,amsthm}
\usepackage{graphicx}
\usepackage{booktabs}
\usepackage{array}
\usepackage{microtype}
\usepackage{placeins}
\usepackage[hidelinks]{hyperref}
\usepackage{cite}
\usepackage{algorithm}
\usepackage{algpseudocode}
\usepackage{authblk}

% IEEE-style numeric citations. The standard ieeetr bibliography style is
% included in common LaTeX installations and requires no custom .bst file.

% Figures are generated in results/figures/main by the scripts in
% experiments/paper_figures. Referencing each figure by filename keeps the
% source portable within the project repository.
\graphicspath{{results/figures/main/}}

\newtheorem{proposition}{Proposition}
\newtheorem{corollary}{Corollary}

\title{Split and Leaf Hijacking:\\
Integrity Attacks on Federated Gradient-Boosted Trees}

\author[1]{Turke Althobaiti}
\author[2]{Habib Ullah Manzoor}
\author[3]{Basim Alhumaily}
\author[2]{Naeem Ramzan\thanks{Corresponding author: \texttt{naeem.ramzan@uws.ac.uk}}}

\affil[1]{Department of Computer Science, Faculty of Science, Northern Border University, Arar 73222, Saudi Arabia \\ \texttt{turke.althobaiti@nbu.edu.sa}}
\affil[2]{School of Computing, Engineering and Physical Sciences, University of the West of Scotland, UK \\ \texttt{habib.manzoor@uws.ac.uk}, \texttt{naeem.ramzan@uws.ac.uk}}
\affil[3]{Department of Electrical Engineering, College of Engineering, Qassim University, Buraydah 52571, Saudi Arabia \\ \texttt{b.alhumaily@qu.edu.sa}}

\date{}

\begin{document}

\maketitle

\begin{abstract}
Organizations such as hospitals and banks often hold tabular data that they cannot share directly. Federated learning (FL) allows them to train a shared model without moving the raw data, and gradient-boosted trees are well suited to this setting. However, federated tree training creates an integrity risk that has received less attention than privacy. At each node, the protocol combines gradient histograms from several parties and selects a split through a discrete argmax. We show that this discrete decision changes how poisoning damage scales. For horizontal FL (HFL), we derive a closed-form split flip margin and prove that it is unchanged when the same rows are repartitioned among clients. Under unrestricted reports, one attacker can reach this margin, so additional colluders do not expand the set of possible aggregate histograms. We introduce split hijacking, which forges or rescales a histogram contribution to change the winning split, and leaf hijacking, which preserves the selected split but corrupts later sample routing. Split hijacking has a sharp threshold, while leaf hijacking causes a gradual response. Experiments on census, clinical, and financial datasets show that one HFL attacker reduces test AUC from $0.9040$ to $0.8953$. When the attacker reaches the full tree, two datasets collapse to AUC $0.50$. In vertical FL (VFL), a label-free backdoor reaches a $12.6\%$ prediction flip rate with a $12.8\%$ AUC cost. These results show that security analysis for federated trees should consider report limits, node coverage, and attack duration, in addition to the number of malicious parties.
\end{abstract}

\section{Introduction}
\label{sec:intro}

Federated gradient-boosted trees are useful for tabular data held by several organizations, including banks and hospitals. Their training protocol differs from the neural network setting studied in most federated learning (FL) security research. Neural FL usually averages updates in a continuous parameter space \cite{mcmahan2017fedavg}. Federated tree boosting instead sums gradient histograms from all clients and selects one feature and one threshold through a discrete argmax at every node of every tree.

This difference changes the poisoning problem. In continuous averaging, an attacker's influence depends on both the number and size of corrupted updates. In a discrete argmax, the immediate goal is to move one candidate past the current winner. Once the candidate crosses that boundary, it wins the split. In our experiments, the selected target remains the winner at every larger injection that we test. Another candidate on the same feature could overtake it at a larger value, but we do not observe this outcome. The main security questions therefore extend beyond the number of malicious parties. They also include whether one report is bounded, how many nodes an attacker can reach, and how long the corruption continues during training.

We study this attack surface in horizontal and vertical federated gradient boosting. In horizontal FL (HFL), each client holds a different set of rows and reports gradient and Hessian sums for every histogram bin. In vertical FL (VFL), each party holds a different set of features, while the active party holds the labels. Passive parties use Paillier's additive homomorphism to build encrypted histograms. These protocols protect data in different ways, but neither protocol proves that a reported histogram was computed honestly.

This paper makes four contributions.

\begin{itemize}
  \item We define a \emph{split flip margin} for HFL. It is the smallest gradient injection into one histogram cell that allows a chosen competitor to match or beat the current winning split. We validate the formula with a full split search at six real internal nodes.

  \item We prove that, for a fixed dataset and shared bins, the HFL margin does not depend on how rows are divided among clients. We also prove that one attacker is sufficient when reports can take any value. This second result applies to any unrestricted additive aggregation, including FedAvg-style weighted sums. The tree-specific result is that reaching a finite, closed-form margin is enough to change a node decision. A per-client bound below that margin restores dependence on the number of colluders.

  \item We introduce two attacks called \emph{split hijacking} and \emph{leaf hijacking}. Split hijacking changes how a split is selected. Leaf hijacking changes how samples are routed after an honest split. Their response curves differ. Split selection changes at a threshold, while the corrupted leaf weight changes continuously with the probability of misrouting.

  \item We evaluate both attacks on census, clinical, and financial datasets. We vary the number of malicious clients, injection size, attack radius, attack duration, federation size, tree depth, target choice, and leaf-misrouting probability. We also build a VFL backdoor that changes selected predictions without inferring labels.
\end{itemize}

Our claims are limited to the evaluated setting. The strongest HFL attack is adaptive and receives the current honest aggregate as additional knowledge. The VFL attacker never sees plaintext gradients. Most experiments use our own protocol-faithful implementations rather than a production system. The default HFL attacker reaches only the root node. We report separate experiments with a larger radius because node coverage, rather than attacker count, causes the most severe failures.

The rest of the paper is organized as follows. Section~\ref{sec:background} reviews federated tree training and related work. Section~\ref{sec:threatmodel} defines the attacker capabilities. Section~\ref{sec:methodology} presents the security parameters, split flip margin, and attack methods. Section~\ref{sec:setup} describes the datasets and experimental design. Section~\ref{sec:experiments} reports the results. Section~\ref{sec:discussion} discusses the implications, limitations, and possible defenses.

\section{Background and Related Work}
\label{sec:background}

\subsection{Gradient boosting and federated histograms}

A gradient-boosted ensemble is built one tree at a time. Let $F_t$ denote the ensemble after $t$ trees, $F_0$ a constant initial prediction, and $\eta$ the learning rate:
\begin{equation}
F_t(x) = F_{t-1}(x) + \eta T_t(x), \qquad
T_t = \arg\min_T \sum_j L\bigl(y_j,F_{t-1}(x_j)+T(x_j)\bigr).
\label{eq:boosting}
\end{equation}
We use a second-order expansion of the loss. This allows tree $T_t$ to be fitted from the gradient $g_j$ and Hessian $h_j$ of each sample under the current ensemble \cite{chen2016xgboost}. Scoring a candidate split requires the sums of gradients and Hessians in each histogram bin.

\paragraph{Horizontal training.}
Each of the $K$ clients holds a different set of rows from the training set $\mathcal{D}$. We use a standard class-conditional Dirichlet split to create label skew. For each class $y\in\mathcal{Y}$:
\begin{equation}
\pi_y \sim \mathrm{Dirichlet}(\alpha\mathbf{1}_K), \qquad
\bigl|\mathcal{D}^{(i)}_y\bigr| \approx \pi_y^{(i)}\lvert\mathcal{D}_y\rvert.
\label{eq:dirichlet}
\end{equation}
Smaller values of $\alpha$ create stronger label skew. Unless stated otherwise, we use $\alpha=0.5$.

At each node, client $i$ reports its local gradient and Hessian sums for feature $f$ and bin $b$. The server adds the reports to form the global histogram:
\begin{equation}
G[f,b] = \sum_{i=1}^K g^{(i)}[f,b], \qquad
H[f,b] = \sum_{i=1}^K h^{(i)}[f,b].
\label{eq:hflagg}
\end{equation}
The server scores every valid feature-threshold pair and selects the pair with the highest gain. Our HFL protocol sums histograms directly. This differs from systems such as SimFL, which find similar instances across parties and train trees one client at a time without creating a shared histogram \cite{li2020fedgbdt}. FederBoost also forms a shared histogram in its horizontal mode, but protects it with a lightweight secure aggregation protocol rather than the unprotected sum our threat model assumes \cite{tian2020federboost}.

\paragraph{Vertical training.}
In VFL, an active party holds the labels and usually some features. Passive parties hold the remaining, disjoint feature blocks \cite{cheng2021secureboost}. The active party computes $g_j$ and $h_j$ for each sample, encrypts them with its Paillier public key, and sends the ciphertexts to the passive parties. Paillier supports ciphertext addition and multiplication of a ciphertext by a plaintext value:
\begin{equation}
\mathrm{Enc}(m_1)\mathrm{Enc}(m_2)=\mathrm{Enc}(m_1+m_2), \qquad
\mathrm{Enc}(m)^s=\mathrm{Enc}(sm).
\label{eq:paillier}
\end{equation}
These operations allow a passive party to build an encrypted histogram without seeing a plaintext gradient:
\begin{equation}
\mathrm{Enc}\bigl(G_p[f,b]\bigr)
= \prod_{j:\,x_j\in\mathrm{bin}(f,b)} \mathrm{Enc}(g_j).
\label{eq:vflagg}
\end{equation}
The active party decrypts the candidate statistics and selects the split. If the winning feature belongs to a passive party, the active party learns the selected feature and threshold as part of the standard SecureBoost protocol. Real systems store signed fixed-point values in $\mathbb{Z}_N$. Our analysis assumes that these values do not wrap around. Section~\ref{sec:setup} confirms that all tested values remain far within the safe range. Other VFL tree systems choose different cryptography over the same Paillier-only design: Pivot combines threshold partially homomorphic encryption with secure multiparty computation to close two plaintext-tree leakages that a SecureBoost-style protocol exposes \cite{wu2020pivot}.

\paragraph{Privacy does not guarantee report integrity.}
Paillier's semantic security hides plaintext gradients from parties without the private key. Our attacks do not break this property. They use only public ciphertext operations that are already available to a passive party. Similarly, HFL histogram summation does not verify that a client's local sums are honest. Privacy and integrity are therefore separate requirements. A protocol can hide a value while still accepting a manipulated report.

\subsection{Related work}

Much of the security research on FL and federated trees focuses on privacy. Gradient inversion reconstructs training inputs from shared gradients \cite{zhu2019deepleakage}. Membership inference tests whether a record was used during training \cite{shokri2017membership}. In tree-based FL, related attacks reconstruct data from histograms, split values, and decision paths. Other attacks infer labels from routing information visible to a passive party \cite{luo2021feature,digennaro2025timberstrike}. Weng et al. use SecureBoost's homomorphic operations in reverse-sum and reverse-multiplication attacks to recover private information \cite{weng2020privacy}. Our VFL attack uses the same cryptographic flexibility for a different purpose. It changes the selected split rather than recovering a secret.

Byzantine-robust aggregation rules such as coordinate-wise median, trimmed mean, and Krum give statistical guarantees against a bounded fraction of corrupted updates \cite{yin2018byzantine,blanchard2017krum}. Later work shows that a small, carefully crafted local update can still poison FedAvg or defeat several of these Byzantine-robust rules \cite{bhagoji2019analyzing,fang2020local}. More recent defenses target colluding clients directly or bootstrap trust from a small clean dataset: FoolsGold flags Sybils from the similarity of their update directions \cite{fung2018foolsgold}, and FLTrust scores each client update against a server-held reference model \cite{cao2021fltrust}. Bagdasaryan et al. show that scaling a malicious update by the inverse of the server's expected aggregation weight lets a single client implant a persistent backdoor in FedAvg \cite{bagdasaryan2020backdoor}. FedClamp identifies anomalous clients from their per-round update statistics, and a related layer-based aggregation defense filters backdoor updates by scoring each layer of a client's update separately \cite{manzoor2022fedclamp,manzoor2023defending}. In applied energy-forecasting settings, false data injection and general data-integrity attacks against federated load forecasting have both been studied, together with matching resilience and robustness evaluations \cite{shabbir2024resilience,shabbir2024robustness}. Other work in this applied setting studies a stealthy attack hidden in the communication round itself, compares adversarial robustness across centralised and decentralised deployments, and proposes an energy- and privacy-aware training scheme for resilience \cite{manzoor2025novel,manzoor2024centralised,manzoor2024enhanced}. A class-coverage sampling and entropy-weighted aggregation scheme addresses client heterogeneity rather than adversarial robustness directly \cite{chen2026class}. A broader survey covers the wider landscape of FL security strategies across models, data, and privacy \cite{manzoor2024survey}. These methods treat each client update as an estimate in a shared continuous parameter space. An HFL histogram is different because each report is a disjoint partial sum and the server needs the exact total. A median or trimmed mean of per-client bin sums would not recover this total. Unequal client sample counts under non-IID splits make this mismatch even larger. Such a defense would alter the honest computation before an attack occurs.

Heterogeneity also affects the two settings differently. In FedAvg, several local nonlinear optimization steps can cause client drift, and robust aggregation becomes harder when client distributions differ. Methods such as bucketing address this problem \cite{karimireddy2022bucketing}. In histogram-based tree boosting, each report is a one-step additive sum. Regrouping the same sample contributions across clients does not change the global total. Section~\ref{sec:marginmethod} states this distinction more precisely.

Label-free backdoors have been studied in vertically split neural networks, where the attacker uses gradients from learned feature embeddings \cite{shen2025labelfree}. Gradient-boosted trees do not contain such embeddings. To the best of our knowledge, the backdoor studied here is the first label-inference-free backdoor for federated tree boosting. It disrupts selected predictions, but we do not show that it reliably moves them to one fixed target class. In continuous-parameter FL, gradient norm clipping and client-level differential privacy noise have both been shown to reduce backdoor success \cite{sun2019backdoor,geyer2017dpfed}. Histogram summation has no direct analogue of a gradient norm to clip, since a bin sum can grow arbitrarily large with the number of local samples routed to it; adapting either defense to tree boosting is an open problem.

\section{Threat Model}
\label{sec:threatmodel}

In HFL, the server follows the aggregation and split-selection protocol. In VFL, the active party also follows the protocol. One or more clients or passive parties may violate the reporting or routing rules. Table~\ref{tab:threatmodel} summarizes the capabilities of each attacker.

\begin{table}[htbp]
\centering
\caption{Attacker capabilities by protocol and mechanism.}
\begin{tabular}{lccl}
\toprule
Attacker & Split hijacking & Leaf hijacking & Knowledge \\
\midrule
Adaptive HFL client & Arbitrary histogram & N/A & Honest aggregate \\
Blind HFL client & Noisy histogram & N/A & Local report only \\
VFL passive party & Ciphertext rescaling & False routing & No plaintext gradients \\
\bottomrule
\end{tabular}
\label{tab:threatmodel}
\end{table}

\paragraph{Adaptive HFL attacker.}
In a normal deployment, a client would not see the other clients' reports before sending its own. Our strongest HFL attack assumes that the attacker knows the current honest aggregate. This is a worst-case assumption, not a proposed side channel. It allows the attacker to identify and strengthen a weak candidate split. We also test a blind baseline that adds Gaussian noise without this knowledge. The baseline tests whether the saturation pattern remains without access to the aggregate, but it does not target the split flip margin directly.

\paragraph{Blind VFL attacker.}
A malicious passive party can apply Paillier's public operations to its ciphertexts and can send false routing bits after a split on one of its features. It never observes a plaintext gradient. For a nonzero plaintext gradient $g$, the party could in principle reach a target value $v$ by choosing $s=v/g$. It cannot compute that scale because it does not know $g$. We therefore test many scale values $s$ over several orders of magnitude. If $g=0$, rescaling cannot change the cell.

The closed-form margin in Section~\ref{sec:marginmethod} applies only to HFL. In HFL, modifying one reported bin also changes the node total reconstructed by the server. In VFL, the active party computes the node total from its labels and keeps it fixed while checking a passive party's histogram. Both score functions are quadratic in the injected gradient, but their coefficients differ.

We do not assume a compromised server, a broken Paillier cryptosystem, fake identities, or help from outside parties. Unless stated otherwise, malicious parties act independently. Each attack uses a normal participant interface by sending a false local statistic, applying an allowed ciphertext operation, or returning a false routing decision.

\section{Methodology}
\label{sec:methodology}

We describe each attack along three security axes, derive the HFL split flip margin, and define the split and leaf hijacking methods. Section~\ref{sec:experiments} evaluates these mechanisms under the setup in Section~\ref{sec:setup}.

\subsection{Security parameters}
\label{sec:parameters}

\paragraph{Byzantine fraction.}
For $K$ participants and a malicious set $M$, define
\begin{equation}
f = \frac{\lvert M\rvert}{K}.
\label{eq:fdef}
\end{equation}
This fraction is a central parameter in research on continuous FL. In our unrestricted HFL threat model, however, the reachable change to the aggregate is already unbounded when $\lvert M\rvert=1$. Section~\ref{sec:splithijack} proves this result and explains why a per-client limit restores dependence on $f$.

\paragraph{Attack radius.}
Let $D$ be the maximum tree depth, with the root at level one. An attacker with radius $r$ corrupts
\begin{equation}
\mathcal{N}(r)=\{v\in T:\mathrm{depth}(v)\leq r\},
\qquad r\in\{1,\ldots,D\}.
\label{eq:rdef}
\end{equation}
Thus, $r=1$ reaches only the root and $r=D$ reaches the full tree. Radius is not the same as the fraction of covered nodes. In a balanced tree, the covered fraction is $(2^r-1)/(2^D-1)$. The ratio $r/D$ describes only the fraction of depth. Our experiments vary $r$ directly. We discuss, but do not prove, whether $r/D$ predicts damage across trees with different depths.

\paragraph{Persistence.}
For $N$ boosting rounds $\mathcal{R}=\{1,\ldots,N\}$ and a set of corrupted rounds $\mathcal{R}_c$, define
\begin{equation}
\tau=\frac{\lvert\mathcal{R}_c\rvert}{\lvert\mathcal{R}\rvert}.
\label{eq:taudef}
\end{equation}
Later trees can partly correct an early corrupted tree because they fit residuals computed after that error. For squared loss, the next residual is $r_{t+1}=y-F_t=r_t-\eta T_t(x)$. A corrupted tree $T_t$ therefore changes the targets used in later rounds. This creates a chance for recovery, but does not guarantee it because a depth-limited tree cannot always undo an arbitrary earlier tree. We measure both short and sustained corruption rather than assuming complete recovery.

\begin{table}[htbp]
\centering
\caption{Security parameters and their roles in the evaluated threat model.}
\begin{tabular}{p{3.0cm}p{4.5cm}p{5.4cm}}
\toprule
Parameter & Common role in continuous FL & Role studied here \\
\midrule
Byzantine fraction $f$ & Primary robustness parameter & Flat after one attacker only when reports are unrestricted \\
Attack radius $r$ & No direct analogue & Controls how much of each tree can be corrupted \\
Persistence $\tau$ & Often fixed or implicit & Separates transient recovery from sustained damage \\
\bottomrule
\end{tabular}
\label{tab:parameters}
\end{table}

\subsection{Horizontal split flip margin}
\label{sec:marginmethod}

At one node, let $(G_L,H_L)$ and $(G_R,H_R)$ denote the left and right gradient and Hessian sums. Let $G=G_L+G_R$ and $H=H_L+H_R$. A standard regularized split score is
\begin{equation}
\mathrm{Gain}
=\frac{1}{2}\left(
\frac{G_L^2}{H_L+\lambda}
+\frac{G_R^2}{H_R+\lambda}
-\frac{G^2}{H+\lambda}
\right)-\gamma.
\label{eq:gain}
\end{equation}
The factor $1/2$ and the constant $-\gamma$ do not affect the winning split at a fixed node. We therefore define
\[
S=\frac{G_L^2}{H_L+\lambda}
+\frac{G_R^2}{H_R+\lambda}
-\frac{G^2}{H+\lambda}.
\]

Fix the current winning score $S_w$ and a target competitor with unmodified score $C$. The \emph{directional one-cell split flip margin} is the smallest $\delta\geq0$ that can be added to the target's left gradient sum so that its score reaches $S_w$. All other reported values remain fixed. This definition does not minimize over every possible perturbation. It excludes negative injections, Hessian changes, and changes to multiple cells or features.

In HFL, adding $\delta$ to a bin in the target's left partition changes $G_L$ to $G_L+\delta$ and the reconstructed node total $G$ to $G+\delta$. The target's right sum does not change. Substitution gives
\begin{equation}
S(\delta)=A\delta^2+B\delta+C,
\qquad
A=\frac{1}{H_L+\lambda}-\frac{1}{H+\lambda},
\qquad
B=\frac{2G_L}{H_L+\lambda}-\frac{2G}{H+\lambda}.
\label{eq:quadratic}
\end{equation}
Every valid candidate satisfies $H_L,H_R\geq h_{\min}>0$. Since $H=H_L+H_R$, we have $A>0$. We assume a strict initial gap, $C<S_w$, which holds when the target is a true runner-up. A tied target has margin zero and is handled by the fixed tie-breaking rule. Under a strict gap, $S(0)-S_w=C-S_w<0$. The parabola opens upward because $A>0$, so its two roots lie on opposite sides of zero. The unique nonnegative crossing is
\begin{equation}
\delta^\star
=\frac{-B+\sqrt{B^2-4A(C-S_w)}}{2A}.
\label{eq:margin}
\end{equation}

Equation~\ref{eq:margin} compares the current winner with one target candidate. Changing one histogram cell can also alter other thresholds on the same feature, so a third candidate may overtake both candidates first. For this reason, our validation reruns the complete split search instead of checking only this pair.

The formula does not directly apply to VFL. In that setting, the active party keeps $G$ fixed and calculates $G_R=G-G_L$. An injection changes $G_L$ and $G_R$ in opposite directions. The score remains a convex quadratic in the injection, so a threshold still exists, but the coefficients differ. We do not derive or use a closed-form VFL margin.

\subsubsection{Properties of the margin}

First, Equation~\ref{eq:margin} should be tight. A perturbation just below it should preserve the current winner, while a perturbation just above it should change the winner unless another candidate crosses first.

Second, the margin should increase with node size. If the data are replicated exactly by a factor $c$ and $\lambda=0$, then $S_c(c\delta)=cS_1(\delta)$, so the margin scales linearly with $c$. When $\lambda>0$, the regularizer does not scale with the data and the measured margin can move away from this linear baseline. The exponent measured across different real nodes is descriptive because node depth, class balance, and the gap between candidates also change with node size.

Third, repartitioning rows does not change the margin when the dataset, bin edges, model state, and tie-breaking rule remain fixed.

\begin{proposition}[Heterogeneity invariance]
\label{prop:heterogeneity}
Fix a training set $\mathcal{D}$, shared bin edges, a tie-breaking rule, and a client count $K$. Consider any two partitions of the rows of $\mathcal{D}$ into $K$ disjoint client sets. Honest HFL builds the same tree node by node under both partitions. Therefore, the split flip margin is identical at every matching node.
\end{proposition}

\begin{proof}
At the root, each aggregate bin is the sum of the same sample gradients and Hessians over all rows in that bin. Addition is associative, so regrouping rows among clients does not change any aggregate bin. The candidate set, scores, selected split, and margin are therefore unchanged. The selected split sends the same samples to each child because routing depends on feature values, not client ownership. The same argument applies to every child by induction. Both runs produce the same predictions after each tree and therefore calculate the same gradients and Hessians in the next round. The full ensemble and every margin remain identical.\qedhere
\end{proof}

This proposition concerns only the aggregate statistic. It does not address detectability or attacker knowledge. A local plausibility test, report-size bound, or detector based on a client's history could still depend on the partition. The proof also assumes shared bin edges. Our simulator builds them with centralized access to all feature values. A deployed system would instead use a federated quantile sketch. Such a sketch aims to estimate the same global boundaries, but we have not tested whether it preserves exact invariance.

The comparison with FedAvg also needs care. A correctly weighted average of fixed sample statistics is invariant to regrouping for the same reason. FedAvg becomes sensitive to heterogeneity because each client reports the result of several local optimization steps rather than one fixed additive statistic. Nonlinear local training and client drift break the simple associativity argument.

\subsection{Attack taxonomy}
\label{sec:attackmethod}

We study two mechanisms. \emph{Split hijacking} changes which candidate is selected by the node's argmax. \emph{Leaf hijacking} preserves that choice but changes which samples reach each child. Both attacks use simple participant operations, but their response to attack strength differs. Figure~\ref{fig:framework} traces both attacks through the federated boosting pipeline, from local histogram contributions to test performance.

\begin{figure}[htbp]
\centering
\includegraphics[width=\textwidth]{framework.pdf}
\caption{Both attacks use the same federated boosting pipeline, which contains local histogram construction, global aggregation, split selection, sample routing, and prediction. Split hijacking acts during aggregation and changes the node's argmax. Leaf hijacking preserves the selected split but corrupts later routing. Split hijacking has a threshold response, while leaf hijacking changes with the misrouted fraction.}
\label{fig:framework}
\end{figure}

\FloatBarrier

\subsubsection{Split hijacking}
\label{sec:splithijack}

An HFL client performs split hijacking by forging one or more histogram cells. A VFL passive party performs it by rescaling an encrypted contribution. Once the target score passes the original winning score, the target remains ahead at all larger injections we test. This does not prove that the final node decision can never change again. Another candidate on the same feature could later overtake the target if its quadratic coefficient were larger. We do not observe this behavior. In the tested range, the winner changes like a step function. We confirm this pattern with the full argmax search rather than assuming it.

When HFL reports are unrestricted, the relationship between attacker count and reachable outcomes is exact.

\begin{proposition}[Single-attacker sufficiency]
\label{prop:singleattacker}
Fix one node and the honest report $h_i\in\mathbb{R}$ of every HFL client at a target cell. Let the honest aggregate be $H_0=\sum_{i=1}^K h_i$. Suppose each malicious client can replace its report with any real value. Every aggregate reachable by $k\geq1$ malicious clients is then reachable by one malicious client, and every one-client aggregate is also reachable by $k$ clients.
\end{proposition}

\begin{proof}
Write a malicious report as $h_i+d_i$, where $d_i$ is unrestricted. For a malicious set $M$, the aggregate becomes
\[
H_0+\sum_{i\in M}d_i.
\]
A single malicious client $a$ can produce the same value by reporting $h_a+\sum_{i\in M}d_i$. Conversely, a joint attack can reproduce any one-client deviation by setting every other malicious deviation to zero. Thus, every nonempty malicious set reaches the same collection of aggregate histograms. Split selection and all later computations are deterministic functions of this aggregate, so the reachable outcomes are also identical.
\end{proof}

\begin{corollary}[Bounded deviations restore count sensitivity]
\label{cor:bounded}
Suppose each malicious deviation is bounded by $\lvert d_i\rvert\leq B$, where $0<B<\delta^\star$. In the required direction, $k$ colluders can produce a joint deviation of at most $kB$. Reaching the margin therefore requires
\[
k\geq\left\lceil\frac{\delta^\star}{B}\right\rceil.
\]
\end{corollary}

Proposition~\ref{prop:singleattacker} does not mean that attacker count is irrelevant for every discrete aggregation rule. It applies to reports that may take any value. Clipping, range proofs, per-record limits, and other verifiable bounds can restore dependence on the Byzantine fraction. We do not evaluate the bounded case.

The proposition is also not unique to trees. Its proof requires only additive aggregation and unrestricted malicious reports. The same reasoning applies to a FedAvg-style weighted sum $\sum_i w_i u_i$. One client can set the sum to any target $A^\star$ by reporting $u_a'=(A^\star-\sum_{i\neq a}w_i u_i)/w_a$. The tree-specific result is narrower: a finite, closed-form node margin is enough to change a choice among a fixed set of candidate splits. If each client's permitted deviation is smaller than this margin, colluders must combine their reports to cross it. Continuous systems can also have margins, robustness radii, and optimization tolerances. Our claim is only that ordinary continuous aggregation does not itself create the same explicit node-local argmax threshold.

The VFL experiment studies a different form of saturation. One passive party owns all candidate splits in its feature block. We vary the size of that party's injection, not the number of colluding passive parties. An unrestricted rescaling can make one candidate dominate a node, but several malicious parties with different feature blocks could affect different parts of the tree. We do not study VFL collusion.

Algorithm~\ref{alg:splithijack} gives the HFL and VFL implementations. The adaptive HFL branch chooses the current worst candidate and assigns it a deliberately large budget. It overshoots the exact margin because a narrow win provides no additional node-level benefit. The VFL branch targets one cell using only information visible to the passive party.

The HFL implementation is stronger than the one-cell perturbation used to define the margin. It zeros the client's full gradient histogram except for the target cell, which receives the full budget. By contrast, the margin analysis adds $\delta$ to one cell and leaves every other cell unchanged. We validate this cleaner primitive independently against aggregate histograms, so the margin check does not depend on the attack implementation.

This distinction also matters for conservation checks. Even the cleaner one-cell attack would fail a cross-feature conservation test because the total for one feature would differ by $\delta$ from the totals for the untouched features. The implementation's zeroing makes the violation more visible, but does not create the underlying problem. We also repeated the main HFL saturation experiment with this cleaner, one-cell construction: it produces the same collapse, from AUC $0.9040$ to $0.8970$ at the first tested attacker fraction (Section~\ref{sec:saturationresults}), closing the gap between the validated primitive and the implemented attack.

\begin{algorithm}[htbp]
\caption{Split hijacking in horizontal and vertical protocols}
\label{alg:splithijack}
\begin{algorithmic}[1]
\Require protocol; target node; regularizers $\lambda,\gamma$; minimum child Hessian $h_{\min}$
\If{protocol is Horizontal}
  \Require local histogram $(g,h)$; honest aggregate of other clients $(g_{-i},h_{-i})$; multiplier $m$
  \For{each candidate $(f,b)$}
    \State Compute $G_L,H_L,G_R,H_R,G,H$ from $(g_{-i}+g,h_{-i}+h)$
    \If{$H_L<h_{\min}$ or $H_R<h_{\min}$}
      \State \textbf{continue}
    \EndIf
    \State Compute $\mathrm{Gain}(f,b)$ using Equation~\ref{eq:gain}
  \EndFor
  \State $(f^\star,b^\star)\gets\arg\min_{(f,b)}\mathrm{Gain}(f,b)$
  \State $\mathrm{budget}\gets m\lVert g_{-i}+g\rVert_1$
  \State $\tilde g\gets\mathbf{0}$; $\tilde g[f^\star,b^\star]\gets\mathrm{budget}$; $\tilde h\gets h$
  \State \Return $(\tilde g,\tilde h)$
\Else
  \Require target feature $f^\star$; encrypted histogram $\mathrm{Enc}(g[f,b])$; scale $s$
  \State $b^\star\gets\arg\max_b\#\{j:x_j\in\mathrm{bin}(f^\star,b)\}$
  \State $\mathrm{Enc}(g[f^\star,b^\star])\gets\mathrm{Enc}(g[f^\star,b^\star])^s$
  \State Leave the Hessian histogram unchanged
  \State \Return the modified encrypted histogram
\EndIf
\end{algorithmic}
\end{algorithm}

\paragraph{Label-inference-free backdoor.}
A malicious passive party can repeatedly force a split on one of its features. At inference time, it can change that feature so that a selected input crosses the planted threshold. Neither step requires a label. For a selected test set $\mathcal{T}$, we report
\begin{equation}
\mathrm{FlipRate}
=\frac{1}{\lvert\mathcal{T}\rvert}
\sum_{x\in\mathcal{T}}
\mathbb{1}\left[\hat y_{\mathrm{backdoor}}(x)
\neq\hat y_{\mathrm{honest}}(x)\right].
\label{eq:fliprate}
\end{equation}
Both models are evaluated on the same triggered input. This isolates the effect of planting the trigger during training from the effect of changing the feature value at test time.

Flip rate is not the success rate of a fixed-target attack. Predictions change in both directions. For example, one run produced 21 flips from class 0 to class 1 and 12 flips in the opposite direction. We therefore describe this mechanism as on-demand prediction disruption rather than a targeted backdoor. We do not measure false triggers or restrict the test set to one source class, so the results do not establish a standard targeted backdoor.

\begin{algorithm}[htbp]
\caption{Label-inference-free backdoor through split hijacking}
\label{alg:backdoor}
\begin{algorithmic}[1]
\Require trigger feature $f_t$; trigger bin $b_t$; scale $s$; plant count $n$
\State $\mathrm{remaining}\gets n$
\For{each tree root that receives this party's histogram}
  \If{$\mathrm{remaining}>0$}
    \State $\mathrm{Enc}(g[f_t,b_t])\gets\mathrm{Enc}(g[f_t,b_t])^s$
    \State $\mathrm{remaining}\gets\mathrm{remaining}-1$
  \EndIf
\EndFor
\Statex
\Require selected input $x$; planted threshold $\theta$
\If{$x[f_t]<\theta$}
  \State Set $x[f_t]$ above the maximum training value
\EndIf
\State \Return $\mathrm{model.predict}(x)$
\end{algorithmic}
\end{algorithm}

\paragraph{Bagging controls.}
To separate discrete split selection from general tree poisoning, we compare boosting with federated bagging. Each client trains a local random forest, and the server averages predicted probabilities:
\[
F_{\mathrm{bag}}(x)=\frac{1}{K}\sum_{k=1}^K f_k(x).
\]
We evaluate label shuffling under both bagging and boosting. We also test a report-level bagging forgery in which a malicious client trains honestly but reports $1-f_k(x)$. Because a probability is bounded to $[0,1]$, this control cannot match the size of the unrestricted histogram perturbation. It compares the shapes of the response curves, not attacks with equal budgets.

\paragraph{Persistence conditions.}
We compare attacks on the first boosting round, the final round, and every round. The transient experiment tests only two positions, both with $\tau=1/20$. It does not provide a complete response curve over all possible schedules.

\subsubsection{Leaf hijacking}
\label{sec:leafhijack}

Leaf hijacking occurs after an honest split on a passive party's feature. The party reports false routing decisions for a selected fraction of the samples that reach the node.

\begin{algorithm}[htbp]
\caption{Leaf misrouting}
\label{alg:leafmisroute}
\begin{algorithmic}[1]
\Require sample indices $I$; selected threshold; misroute probability $p$
\For{each $i\in I$}
  \State $t_i\gets\mathbb{1}[x_i\leq\mathrm{threshold}]$
  \State Draw $u_i\sim\mathrm{Uniform}(0,1)$
  \If{$u_i<p$}
    \State Report $\lnot t_i$
  \Else
    \State Report $t_i$
  \EndIf
\EndFor
\end{algorithmic}
\end{algorithm}

\begin{proposition}[Continuity of the corrupted leaf weight]
\label{prop:leafcontinuity}
Consider a selected split with true partition sums $(G_L,H_L)$ and $(G_R,H_R)$. Suppose each sample is independently misrouted with probability $p$. The expected fitted weight of the corrupted left child is then a continuous function of $p$ on $[0,1]$.
\end{proposition}

\begin{proof}
The attack changes routing but does not alter an individual gradient or Hessian. Therefore,
\begin{equation}
\mathbb{E}[G_L(p)]=(1-p)G_L+pG_R, \qquad
\mathbb{E}[H_L(p)]=(1-p)H_L+pH_R.
\label{eq:leafcont}
\end{equation}
For an outcome $z\in\{0,1\}^n$ of the routing flips, let $w(z)=-G_L(z)/(H_L(z)+\lambda)$. Since $\lambda>0$ and the Hessians are nonnegative, every $w(z)$ is finite. The expected value is
\begin{equation}
\mathbb{E}[w(p)]
=\sum_{z\in\{0,1\}^n}
w(z)p^{\lVert z\rVert_1}(1-p)^{n-\lVert z\rVert_1}.
\label{eq:bernstein}
\end{equation}
This finite Bernstein polynomial is continuous on $[0,1]$ and is infinitely differentiable on the interval.
\end{proof}

The proposition concerns one leaf weight, not test AUC. AUC is a rank statistic on a finite test set and can change in discrete steps even when prediction scores move smoothly. Misrouting at several nodes can also compound. The proposition explains why we expect the smooth model-level curves in Section~\ref{sec:leafresults}, but it does not prove those curves.

\paragraph{Combined attack.}
We combine a small ciphertext rescaling with a small misrouting probability. Let the damage of condition $X$ be $\Delta_X=\mathrm{AUC}_{\mathrm{honest}}-\mathrm{AUC}_X$. We define the interaction statistic
\begin{equation}
\mathrm{Syn}
=\Delta_{\mathrm{split}+\mathrm{leaf}}
-\left(\Delta_{\mathrm{split}}+\Delta_{\mathrm{leaf}}\right).
\label{eq:synergy}
\end{equation}
A positive value indicates more-than-additive damage. Zero indicates additive damage, and a negative value indicates that the attacks partly mask each other. We estimate this statistic with paired seeds rather than separate mean curves.

\section{Experimental Setup}
\label{sec:setup}

\subsection{Datasets and preprocessing}

We evaluate three binary classification datasets from different domains: UCI Adult, the Cleveland subset of UCI Heart Disease, and UCI Default of Credit Card Clients. We remove rows with missing values and encode categorical variables as integers. The positive class represents income above \$50K, the presence of heart disease, or default in the following month. Table~\ref{tab:datasets} summarizes the role of each dataset.

\begin{table}[htbp]
\centering
\caption{Datasets and the rationale for each VFL feature partition. Row counts are approximate after preprocessing.}
\begin{tabular}{lrrp{7.0cm}}
\toprule
Dataset & Rows & Domain & VFL partition \\
\midrule
Adult & $32{,}000$ & Census & Demographics versus employment and financial attributes \\
Heart Disease & $297$ & Clinical & Demographics and history versus clinical test results \\
Credit Default & $30{,}000$ & Financial & Demographics and credit limit versus payment and billing history \\
\bottomrule
\end{tabular}
\label{tab:datasets}
\end{table}

Adult is used for the main parameter sweeps. The other datasets test whether the findings transfer across sample sizes, class balances, and feature meanings. For VFL, we choose feature blocks that resemble realistic divisions between institutions rather than arbitrary column splits.

\subsection{Training configuration}

The default HFL setup uses five clients, 32 histogram bins, a maximum depth of four, and 20 boosting rounds. Unless a sweep changes it, the Dirichlet concentration is $\alpha=0.5$. The VFL setup uses one active party, one passive party, 16 bins, and a maximum depth of three. Homomorphic computation is expensive, so each VFL run first draws at most 1,500 rows and then divides them into training and test sets. Adult and Credit Default therefore use 1,500-row VFL subsamples, while Heart Disease uses all 297 rows.

Both protocols use a learning rate of $\eta=0.3$, an $L_2$ regularizer of $\lambda=1.0$, and a split penalty of $\gamma=0$. A child is valid only when $H_L,H_R\geq h_{\min}=1.0$. HFL reserves $20\%$ of the data for testing, and VFL reserves $25\%$. The HFL split-hijacking multiplier is $m=3$, so the forged cell receives three times the $L_1$ norm of the honest aggregate gradient histogram.

Both training systems are protocol-faithful implementations written from scratch. The HFL system sums histograms directly. The VFL system implements the published SecureBoost computation \cite{cheng2021secureboost}. We also built and ran the bundled horizontal and vertical examples from an independent reference implementation. Both honest setups trained to convergence and produced baseline behavior consistent with our systems. We did not port the attacks to that codebase, so every attack result reported below applies to our implementations rather than to a named production release.

\subsection{Paillier encoding and arithmetic range}

We use 512-bit Paillier keys to reduce experimental cost. This key size is not secure enough for deployment, but key length does not change the public malleability used by the attack. We also verify that ciphertext rescaling cannot cause wraparound at any tested value.

The \texttt{phe} library represents $x$ as $\lfloor x16^{-e}\rfloor$ for a base-16 exponent $e$. Consider a worst-case bin containing 1,500 gradients, each with $\lvert g\rvert\leq1$, under the largest tested scale of 200. The plaintext magnitude is at most $3\times10^5$. With $e=-9$, it is encoded as the integer $2.06\times10^{16}$. The 512-bit key has an overflow point of $\lfloor n/3\rfloor=3.75\times10^{153}$, which provides about 137 orders of magnitude of headroom. Round-trip tests remain exact and correctly signed for multipliers up to $10^{18}$, far above the values used in the experiments.

\subsection{Statistical analysis}
\label{sec:statmethod}

Unless stated otherwise, each point reports the mean and one standard deviation over five seeds. Each seed redraws the HFL partition. For VFL, it also redraws the data subsample and train-test split. Five seeds provide limited precision, especially for Heart Disease, whose test set contains only about 60 rows.

For the interaction statistic in Equation~\ref{eq:synergy}, each seed uses the same partition and train-test split for the split-only, leaf-only, and combined conditions. We calculate
\begin{equation}
\mathrm{Syn}_i
=\Delta_{\mathrm{split}+\mathrm{leaf}}^{(i)}
-\left(\Delta_{\mathrm{split}}^{(i)}+\Delta_{\mathrm{leaf}}^{(i)}\right),
\qquad
\overline{\mathrm{Syn}}
\mathbin{\pm}
t_{n-1,0.975}\frac{s_{\mathrm{Syn}}}{\sqrt{n}},
\label{eq:synergyci}
\end{equation}
with $n=5$. Pairing avoids treating the shared seed-level data split as independent across conditions.

\subsection{Honest baseline quality}
\label{sec:baselinequality}

Figure~\ref{fig:baseline} compares train and test AUC over boosting rounds for every protocol and dataset. We extend the VFL curves to 30 rounds to check for underfitting and overfitting. Across the six panels, the final test AUC is between 0.0000 and 0.0134 below the peak. This small difference is consistent with seed-level variation rather than clear overfitting.

The Heart Disease VFL curve shows why results should be averaged across seeds. Individual seeds reach their peak at rounds 3, 9, 15, 25, and 3. Round five gives $0.8526\pm0.0126$, while round 30 gives $0.8504\pm0.0199$. We do not perform a formal equivalence test, so the overlap indicates only that there is no clear evidence of a difference. Credit Default VFL improves from $0.7517\pm0.0509$ at round three to a peak of 0.7761 at round eight, which suggests mild undertraining. However, the increase remains within one standard deviation of the round-three estimate.

The reported honest baseline differs slightly between figures, for example $0.9040$, $0.8965$, and $0.9086$ for HFL Adult. These differences come from different experiment scripts drawing different train-test splits or data subsamples under otherwise similar settings, not from different models or protocols. Every results subsection therefore reports the honest baseline for its own configuration next to the attacked value, so each comparison is read against its own reference.

\begin{figure}
\centering
\includegraphics[width=0.8\textwidth]{fig0_baseline_convergence.pdf}
\caption{Train and test AUC over boosting rounds for the honest baseline. Each curve shows the mean and one standard deviation over five seeds. Across the six panels, the difference between peak and final test AUC ranges from 0.0000 to 0.0134.}
\label{fig:baseline}
\end{figure}

\FloatBarrier

\section{Results}
\label{sec:experiments}

We present the results from mechanism to impact. We first validate the split flip margin at individual nodes. We then test saturation over attacker count and injection size. Next, we measure the effects of attack radius, persistence, and federation size. We evaluate the label-free backdoor and compare split hijacking with leaf hijacking. Finally, we examine combined attacks, transfer across datasets, and blind target selection.

\subsection{Split flip margin validation}
\label{sec:marginresults}

Figure~\ref{fig:margin} tests the three margin properties from Section~\ref{sec:marginmethod}. We evaluate six internal nodes from one honest tree. An injection just below Equation~\ref{eq:margin} never changes the full argmax, while an injection just above it always does. This is an exact check at the tested nodes rather than a broad statistical test.

We also test up to 12 nodes at each of six Dirichlet concentrations. The margin grows with node sample count at an estimated power of about 1.8. This is steeper than the exact linear scaling derived for proportional replication when $\lambda=0$. The observations come from different real nodes, so node size changes together with depth, class balance, and the gap between candidates. The exponent should therefore be read as a description, not a controlled scaling law.

The median margin is 90.73 at every tested Dirichlet concentration. This exact match confirms the partition invariance in Proposition~\ref{prop:heterogeneity} for our simulator and its shared global binning.

\begin{figure}[htbp]
\centering
\includegraphics[width=\textwidth]{fig3_margin.pdf}
\caption{Validation of the HFL split flip margin. The closed-form value matches the full split search at six tested nodes. Across naturally different nodes, the margin grows faster than linearly with sample count. It is identical at every tested Dirichlet concentration.}
\label{fig:margin}
\end{figure}

\FloatBarrier

\subsection{Split hijacking: attacker count and injection magnitude}
\label{sec:saturationresults}

Figure~\ref{fig:saturation} compares saturation in HFL, VFL, bagging, and boosting. In HFL, the number of malicious clients ranges from zero to three out of five, giving $f\in\{0,0.2,0.4,0.6\}$. The target histogram attack causes its full AUC loss at $f=0.2$, the smallest nonzero value tested. Additional malicious clients cause no further damage. This result agrees with Proposition~\ref{prop:singleattacker} under unrestricted reports.

Panel (a) also reports a clean one-cell injection, which adds only the closed-form margin to a single histogram cell and leaves every other cell -- including the client's own other features -- untouched, rather than zeroing the rest of the report as the implemented attack above does. This minimal construction produces the same collapse, from AUC $0.9040$ to $0.8970$ at $f=0.2$, and remains exactly at $0.8970$ through $f=0.6$. The saturating collapse therefore holds under the literal primitive Section~\ref{sec:marginmethod}'s margin models, not only under the stronger, more conspicuous construction used elsewhere in this paper.

The blind Gaussian-noise baseline causes less damage but has the same flattening shape. This shows that access to the honest aggregate is not the only cause of saturation. It does not show that a blind attacker can match the adaptive attack in size or precision.

The VFL curve varies one passive party's scale over two orders of magnitude. A moderate scale causes a sharp drop, while larger scales add little damage. This is saturation in injection size, not VFL colluder count. After the passive candidate passes the original winner, the selected split remains fixed at all larger tested scales. A different candidate on the same feature could overtake it at a still larger scale, but we do not observe this outcome.

\paragraph{Aggregation controls.}
We keep label shuffling fixed and compare bagging with boosting. Bagging starts at an honest AUC of 0.8893. Its AUC decreases by 0.0027, 0.0041, and 0.0165 with one, two, and three malicious clients. Boosting starts at 0.9040 and decreases by 0.0033, 0.0035, and 0.0109. The second boosting attacker adds almost no damage, while the third adds more. This control shows flattening but not complete saturation.

A report-level bagging attack makes the contrast clearer. Malicious clients train honestly but report the complement of their predicted probability. From an honest AUC of 0.8910, the losses are 0.0111, 0.0674, and 0.5122 for one, two, and three attackers. The last point has high variation of $\pm0.2223$ because three inverted reports outvote two honest reports in a five-client average. Under averaging, the forgery therefore becomes stronger as attacker count increases.

The comparison is not budget matched. Bagging probabilities are limited to $[0,1]$, while the HFL attacker can submit any real histogram value. These controls compare response shapes in the natural report space of each method. They do not prove that argmax aggregation would fully saturate under the same normalized per-client bound as averaging.

\begin{figure}[htbp]
\centering
\includegraphics[width=0.8\textwidth]{fig1_saturation.pdf}
\caption{Saturation under discrete split selection. Panel (a) varies the number of HFL colluders, and additionally shows a clean one-cell margin injection alongside the implemented attack. Panel (b) varies the scale used by one VFL passive party and is not a collusion experiment. Panels (c) and (d) keep label shuffling fixed and compare bagging with boosting.}
\label{fig:saturation}
\end{figure}

\paragraph{Cross-feature conservation.}
The default HFL forgery fails a simple conservation test because a client's bin sums should produce the same total for every feature. We therefore test a conservation-preserving version that distributes a compensating amount evenly across other features. Its AUC is $0.8968\pm0.0002$, compared with $0.8953\pm0.0019$ for the default attack and 0.9040 for the honest model. However, the conserved attack loses target control. In all five seeds, another inflated feature wins before the intended target. In this construction, conservation checking prevents precise targeting but not untargeted damage. A better tuned conservation-preserving attack might recover target control, but we do not test one.

\paragraph{Sensitivity to VFL training duration.}
Homomorphic computation limits the default VFL training duration, and Section~\ref{sec:baselinequality} shows that the default Adult VFL model is undertrained. We repeat the main VFL attack with longer training closer to convergence. The saturation threshold remains at the same scale, but the damage is smaller. Under the default duration, AUC decreases from 0.8690 to 0.7926, a loss of 0.0764. Under longer training, it decreases from 0.8953 to 0.8806, a loss of 0.0147. The threshold behavior remains, but the default damage should be viewed as an upper bound for an undertrained model.

\FloatBarrier

\subsection{Attack radius, federation size, and tree depth}
\label{sec:radiusresults}

The earlier HFL experiments use a conservative radius of $r=1$ and attack only the root. Figure~\ref{fig:attackdepth} keeps the malicious-client count at one and extends the radius toward the full tree. Damage increases sharply with coverage. On Adult, AUC falls from 0.8965 under a root-only attack to exactly 0.5000 at full depth. On Credit Default, it falls from 0.7365 to exactly 0.5000. Both full-depth results have zero variance over five seeds.

Direct inspection shows that these results are not caused by random predictions. Every test sample receives the same probability: 1.0 on Adult and a value near zero on Credit Default. The models have collapsed to constant predictors. Heart Disease behaves differently, falling from an honest 0.8694 to $0.4546\pm0.0438$ with nonzero variation. Because its test set is much smaller, we do not describe this case as a constant collapse.

\begin{figure}[htbp]
\centering
\includegraphics[width=0.9\textwidth]{fig12_attack_depth_scoping.pdf}
\caption{AUC with one malicious HFL client as the attack radius grows. Full-tree coverage collapses the Adult and Credit Default models to constant predictors with AUC 0.50. Heart Disease remains more variable because its test set is small.}
\label{fig:attackdepth}
\end{figure}

Figure~\ref{fig:federationscaling} checks whether the default federation size explains the saturation result. We keep one malicious HFL client and increase the total client count from 5 to 40. The honest baseline stays at 0.9040 at every client count, and the AUC loss remains between 0.0060 and 0.0087. More honest clients do not dilute the attack because the malicious report remains unrestricted.

For VFL, we divide the same feature space among one, two, or three passive parties while keeping the malicious scale at ten. Honest and attacked AUC agree to four decimal places in all three settings, at 0.8690 and 0.7926 respectively. The attack depends on the target feature block, not on how the other honest features are divided. This experiment does not change the number of malicious passive parties.

\begin{figure}[htbp]
\centering
\includegraphics[width=0.9\textwidth]{fig14_federation_size_scaling.pdf}
\caption{Federation-size checks over five seeds. Panel (a) varies the total number of HFL clients while keeping one attacker. Panel (b) divides the honest VFL feature space among more passive parties while keeping one attacker.}
\label{fig:federationscaling}
\end{figure}

\paragraph{Tree-depth sensitivity.}
Production gradient boosting often uses depths from six to ten. We keep the attack radius at $r=1$ and increase maximum depth. At depth four, the AUC loss is about 0.0087 to 0.0096. It decreases to 0.0011 to 0.0018 at depth six and approaches zero at depths eight and ten. This pattern is consistent with the root becoming a smaller part of the tree's capacity. It does not prove that damage is a function of $r/D$. The radius experiment changes $r$ at fixed $D$, while the depth experiment changes $D$ at fixed $r$. A matched experiment that varies both is needed to support a general scaling law.

\FloatBarrier

\subsection{Persistence across boosting rounds}
\label{sec:persistenceresults}

Figure~\ref{fig:selfheal} compares a corrupted first round, a corrupted final round, and an unattacked baseline on every dataset. Final AUC remains within one standard deviation of the baseline in all transient cases. Adult finishes at 0.9086 without attack, 0.9082 after early corruption, and 0.9086 after late corruption. Heart Disease finishes at 0.8855, 0.8804, and 0.8844. Credit Default finishes at 0.7424, 0.7466, and 0.7418.

The two transient cases recover for different reasons. Early corruption causes an immediate loss, after which several honest trees fit the altered residuals and correct much of the error. Late corruption has little visible effect because the final tree fits a small residual and has limited influence on the converged ensemble. The first case reflects correction, while the second reflects low leverage. Sustained corruption receives neither protection and causes the lasting damage shown in Figure~\ref{fig:saturation}.

These experiments test $\tau=1/20$ at two positions and $\tau=1$. They do not establish a continuous relationship for periodic, intermittent, or middle-round schedules.

\begin{figure}[htbp]
\centering
\includegraphics[width=\textwidth]{fig2_selfheal.pdf}
\caption{Final AUC after corruption of the first or final round, compared with clean training. Each point shows the mean and one standard deviation over five seeds. Sustained corruption appears separately in Figure~\ref{fig:saturation}.}
\label{fig:selfheal}
\end{figure}

\FloatBarrier

\subsection{Label-inference-free backdoor}
\label{sec:backdoorresults}

Figure~\ref{fig:backdoortradeoff} varies the number of trees that contain the same VFL trigger split. With 1, 2, 3, 5, 8, and 13 plants, the AUC costs are $0.8\%$, $0.5\%$, $0.8\%$, $3.2\%$, $12.8\%$, and $30.4\%$. The corresponding flip rates are $9.8\%$, $5.1\%$, $4.1\%$, $8.2\%$, $12.6\%$, and $12.6\%$.

The early differences in flip rate are smaller than their uncertainty. One plant gives $9.8\pm5.0\%$, while two plants give $5.1\pm3.6\%$. We interpret these early differences as noise rather than evidence that additional plants reduce flip rate. The later ceiling is clearer. Eight and thirteen plants give the same flip rate, but thirteen plants more than doubles the AUC cost. By eight plants, the selected samples pass through every relevant planted node. Later plants affect other paths and increase damage without improving the measured trigger effect.

\begin{figure}[htbp]
\centering
\includegraphics[width=0.9\textwidth]{fig9_backdoor_tradeoff.pdf}
\caption{Backdoor tradeoff over five seeds. The AUC cost increases with plant count. Flip rate is noisy and reaches a measured ceiling of $12.6\%$ at eight plants.}
\label{fig:backdoortradeoff}
\end{figure}

\paragraph{Comparison with label inference.}
A standard instance-clustering method that receives true labels for one fifth of the training set reaches $81.3\pm3.1\%$ label-inference accuracy, with individual results from about $78\%$ to $85\%$. Confidence filtering improves accuracy but reduces coverage. A minimum agreement of $70\%$ gives $84.9\%$ accuracy at $88.9\%$ coverage. A minimum agreement of $85\%$ gives $88.7\%$ accuracy at $74.8\%$ coverage. Unanimous voting gives $92.3\pm3.1\%$ accuracy at only $39\%$ coverage.

Label-inference accuracy and backdoor flip rate measure different outcomes. Inference accuracy is measured over candidate labels, while flip rate is measured on selected triggered inputs. The comparison supports only one advantage: the label-free mechanism cannot inherit label-inference errors because it performs no inference. It does not make the backdoor reliable, and the largest measured flip rate remains modest.

\paragraph{Transfer to other datasets.}
On Heart Disease, flip rate increases from $4.7\%$ with one plant to $6.8\%$ with five plants and $11.3\%$ with eight plants, with no clear AUC cost. The eight-plant result is highly uncertain at $11.3\pm11.2\%$. The sweep does not include two and three plants, where Adult shows an early dip, so the two response shapes cannot be distinguished with confidence.

On Credit Default, all three plant counts give a flip rate of $1.6\pm2.1\%$ with almost no AUC cost. Its default VFL setting has only three boosting rounds, so additional plants cannot reach new nodes on the target paths. The backdoor mechanism transfers, but its reliability and cost depend on training duration and path coverage.

\FloatBarrier

\subsection{Leaf hijacking response curve}
\label{sec:leafresults}

Figure~\ref{fig:leafmisdirection} varies misrouting probability from zero to one. As Proposition~\ref{prop:leafcontinuity} suggests, damage increases gradually with attack strength rather than showing the sharp threshold seen in split hijacking. On Adult, AUC decreases from 0.8444 to 0.8393 and 0.8376 at the two smallest nonzero probabilities, then to 0.7690 at half misrouting and $0.2383\pm0.0767$ at full misrouting. Heart Disease decreases from 0.8526 to $0.1570\pm0.0468$ at full misrouting. Credit Default decreases from 0.7194 to $0.3639\pm0.0314$.

The end of the Credit Default curve is not monotone. Its AUC at full misrouting is 0.3639, which is higher than the value of 0.3176 at three-quarter misrouting. Full misrouting reverses the routing rule in a fixed and deterministic way, which can preserve more structure than partial random routing. Heart Disease also has high uncertainty at half misrouting, with $0.4683\pm0.2372$, because its test set is small. The main finding remains unchanged: model damage varies with the misrouting probability and does not reach its maximum at the first nonzero value.

\begin{figure}[htbp]
\centering
\includegraphics[width=\textwidth]{fig15_leaf_misdirection.pdf}
\caption{AUC as leaf-misrouting probability increases from zero to one. Each curve shows the mean and one standard deviation over five seeds on the three datasets.}
\label{fig:leafmisdirection}
\end{figure}

\FloatBarrier

\subsection{Combined split and leaf attacks}
\label{sec:combinedresults}

The largest tested combination of ciphertext rescaling and leaf misrouting decreases AUC from 0.8695 to $0.7173\pm0.0407$, so the effects accumulate. However, we find no clear evidence of a nonadditive interaction. Across four combined settings, the paired synergy estimates and their $95\%$ confidence intervals are
\[
0.003\;[-0.014,0.020],\quad
-0.011\;[-0.048,0.026],\quad
-0.004\;[-0.073,0.065],\quad
0.027\;[-0.043,0.096].
\]
Every interval includes zero. With five seeds and wide intervals, the supported conclusion is that we do not detect synergy. This is not evidence that the attacks are exactly additive. An equivalence test with a defined tolerance and more seeds would be needed to support that claim.

\subsection{Cross-dataset transfer and blind target selection}
\label{sec:crossdataset}

Earlier sections report several findings on all three datasets. Here, we focus on transfer under the default root-only HFL attack and fixed-feature VFL attack, then explain why Credit Default initially appears resistant.

Figure~\ref{fig:crossdataset} applies both default attacks without retuning. Heart Disease reproduces HFL saturation, falling from an honest 0.8694 to 0.8302 at one malicious client and staying flat as more join, and shows a larger VFL loss than Adult, from 0.8521 to 0.6537 at the largest tested scale. Its small sample size leaves less alternative signal after a forced split takes control. Credit Default starts at 0.7423 under HFL and 0.7162 under VFL.

At the default settings, Credit Default initially appears resistant to both attacks for different reasons. In HFL, adaptive target selection works, but root-only coverage is too narrow. The radius sweep in Figure~\ref{fig:attackdepth} shows that wider coverage still collapses the model. In VFL, the blind attacker uses a fixed feature that is a poor target.

We therefore test a label-free VFL heuristic that selects the passive feature with the largest variance in its locally binned values. Figure~\ref{fig:variancetargeting} shows that this choice increases the Credit Default AUC loss from $0.0036\pm0.0184$ to $0.1976\pm0.0990$. On Adult, the loss increases from $0.0764\pm0.0230$ to $0.3512\pm0.0387$. On Heart Disease, the fixed and variance-based losses are $0.1061\pm0.1334$ and $0.1226\pm0.1611$, which overlap widely. We do not perform formal equivalence tests.

The heuristic uses no labels, gradients, or information from another party. Local variance is only a weak indicator of split influence, but it is enough to improve an arbitrary fixed target in two datasets. Finding a stronger blind target-selection rule remains open.

\begin{figure}[htbp]
\centering
\includegraphics[width=0.9\textwidth]{fig7_cross_dataset.pdf}
\caption{Cross-dataset evaluation with the default root-only HFL radius and fixed VFL target feature. Heart Disease reproduces both effects. Credit Default requires wider HFL coverage or a better blind VFL target.}
\label{fig:crossdataset}
\end{figure}

\begin{figure}[htbp]
\centering
\includegraphics[width=0.9\textwidth]{fig13_variance_targeting.pdf}
\caption{Fixed-feature and local-variance VFL target selection over five seeds. The variance heuristic increases damage on Adult and Credit Default. Heart Disease is too variable to distinguish the choices.}
\label{fig:variancetargeting}
\end{figure}

\FloatBarrier

\subsection{Summary of empirical findings}

Table~\ref{tab:summary} separates two effects that are often described with the same word, saturation. In HFL, we test saturation in malicious-client count. In VFL, we test saturation in the injection size of one passive party.

\begin{table}
\centering
\caption{Summary of the findings and the scope of each claim.}
\begin{tabular}{lccp{5.8cm}}
\toprule
Finding & HFL & VFL & Scope \\
\midrule
Attacker-count saturation & Yes & Not tested & HFL count is flat under unrestricted reports \\
Injection-magnitude saturation & N/A & Yes & VFL scale is flat after the selected split changes \\
Transient corruption recovery & Yes & N/A & First and final rounds tested at $\tau=1/20$ \\
Sustained corruption & Yes & N/A & An every-round attack prevents recovery \\
Prediction-disruption backdoor & N/A & Yes & Label-free, but not a fixed-target backdoor \\
Leaf hijacking & N/A & Yes & Damage varies with misrouting probability \\
Cross-dataset transfer & Yes & Yes & Requires sufficient radius or a useful blind target \\
\bottomrule
\end{tabular}
\label{tab:summary}
\end{table}

\FloatBarrier

\section{Discussion and Limitations}
\label{sec:discussion}

\subsection{What the results establish}

The main result is a difference in how damage scales. Under unrestricted HFL reports, additional attackers after the first do not expand the set of reachable aggregate histograms. In contrast, greater node coverage clearly expands the damage. Boosting dynamics create a second distinction. A short attack can be corrected or have little influence, while sustained corruption accumulates. A security evaluation based only on the Byzantine fraction can therefore miss the parameters that control damage in this setting.

The Byzantine fraction is still important when reports are bounded. Corollary~\ref{cor:bounded} shows that a verifiable per-client limit below the split margin forces collusion. This suggests defenses such as range proofs, per-record sensitivity bounds, or proofs that each histogram was computed from committed samples. Robust aggregation rules that estimate a consensus update do not directly solve the problem because honest HFL requires an exact sum of disjoint contributions, not a consensus estimate.

\subsection{Detectability and structural checks}

The evaluated attacks favor clear causal behavior over stealth. The HFL attack injects a large multiple of the honest aggregate, and the VFL sweep uses scales as large as 200. Simple sanity checks may detect these values. For example, the active VFL party can compare a passive feature's reported total with the node total computed from the labels. In HFL, each client's bin sums should produce the same total across all features because every local sample contributes once to each feature histogram.

The conservation-preserving experiment in Section~\ref{sec:saturationresults} shows that this check prevents precise target control in the tested construction but does not prevent untargeted damage. We do not know whether a more carefully tuned conservation-preserving attack could recover target control. A broader defense would require parties to prove that their reports match a committed computation without revealing raw data \cite{xu2020verifynet}. Monitoring reports across rounds may also help because persistence can be visible over time even when one report is plausible under non-IID sampling. Formal detection limits and full defense evaluation remain future work.

These checks all assume the server or active party can inspect a per-client report. A secure aggregation protocol that sums client contributions without revealing them individually \cite{bonawitz2017secureagg} would remove that option entirely, trading detectability for confidentiality. Our threat model assumes the server sees each client's report and does not need to resolve this tension, but a deployment that adds secure aggregation for privacy would also need a different, aggregate-only integrity check.

\subsection{Implementation scope}

The experiments use our own implementations of the two protocols. We compare only honest baseline behavior with a reference system. Deployed systems may include clipping, authentication, secure aggregation \cite{bonawitz2017secureagg}, quantile sketches, auditing, or other checks that change which attacks are possible. We do not provide source code, exact seeds, or the reference system version with this manuscript, which limits independent reproduction.

The HFL simulator uses centralized access to all feature values when constructing shared bin edges. This choice isolates histogram aggregation but does not preserve data privacy. A deployed system would use a federated quantile sketch. Proposition~\ref{prop:heterogeneity} should remain approximately valid if the sketch estimates the same global boundaries, but exact invariance has not been tested under this replacement.

The 512-bit Paillier keys are intentionally too small for deployment. They support the required homomorphic operations and remain far from arithmetic overflow, but they are not a security recommendation.

\subsection{Coverage and tree depth}

Most HFL comparisons attack only the root of a depth-four tree. This design keeps different conditions comparable without collapsing every model to a floor value, but it understates the largest possible effect. The radius sweep shows that one attacker can cause complete failure when it can corrupt the full tree.

The separate radius and depth sweeps in Section~\ref{sec:radiusresults} do not prove that damage follows a general function of $r/D$. A matched design would need to test the same ratios across several depths. Under a root-only attack, the client forges the target node's histogram and reports honestly below that node. The altered root changes the partition, but the attacker does not directly forge later leaf weights. Full-radius experiments repeat the split attack at more nodes and therefore represent a stronger attacker.

\subsection{Vertical undertraining and target selection}

The longer-duration experiment in Section~\ref{sec:saturationresults} shows that VFL saturation remains after additional training, although the damage becomes smaller. The default effect size should therefore be treated as an upper bound for an undertrained model, not as a convergence-independent estimate.

Blind VFL target selection is another limitation. The local-variance rule is effective on two datasets but is not guaranteed to find an influential feature. Stronger label-free rules based on local data structure may improve targeting. Alternatively, feature commitments and range proofs may limit the attacker's choices.

\subsection{Backdoor interpretation}

The VFL backdoor changes predictions in both directions and reaches a maximum measured flip rate of $12.6\%$. It should not be treated as a standard fixed-target backdoor. A stronger evaluation would select inputs from one honest source class, test whether they move to a chosen target class, and measure false triggers on unselected inputs. Our results establish label-free on-demand disruption through one manipulated feature, not reliable control of the output class.

\section{Conclusion}
\label{sec:conclusion}

Federated gradient-boosted trees make repeated discrete decisions from aggregated histograms. Under unrestricted HFL reports, one malicious participant can reproduce any aggregate change that several colluders could create, so additional attackers do not expand the reachable aggregate histograms. In our default experiment, one attacker among five reduces test AUC from $0.9040$ to $0.8953$, with no further loss as more colluders join. The tree-specific result is that success is controlled by a finite, closed-form split flip margin. Across naturally different nodes, this margin scales empirically as $n^{1.77}$ with sample count, although this is a descriptive relationship rather than a controlled law. A verifiable per-client limit below the margin restores dependence on colluder count. The experiments also show that node coverage and persistence are central security parameters. A root-only attack causes a small but repeatable loss, while full-tree coverage collapses Adult from AUC $0.8965$ to $0.5000$ and Credit Default from $0.7365$ to $0.5000$. One corrupted round is corrected or has low influence. On Adult, the final AUC remains within $0.0004$ of the clean baseline when either the first or final round is attacked. Corruption in every round does not recover. In VFL, ciphertext rescaling shows a similar threshold in injection size and lowers AUC from $0.8690$ to $0.7926$ in the default undertrained setting. The same mechanism supports label-free prediction disruption with a $12.6\%$ flip rate at a $12.8\%$ AUC cost. Leaf hijacking behaves differently because it attacks routing instead of the argmax. On Adult, AUC changes gradually from $0.8444$ to $0.2383$ as misrouting increases from zero to one. Overall, security analysis should follow the aggregation and decision rules of the learning protocol. For federated trees, useful guarantees must consider split margins, verifiable report bounds, node coverage, and attack duration. The Byzantine fraction alone does not describe this attack surface.

\section*{Acknowledgements}

During preparation of this manuscript, the authors used an AI-based language model to help edit the prose, verify mathematical derivations, check experimental code, and compare claims with experimental results. The authors reviewed and verified the complete manuscript and take full responsibility for its content, analysis, and conclusions.

\section*{Conflict of Interest}

The authors declare no conflict of interest.

\bibliographystyle{ieeetr}
\bibliography{ref}

\end{document}